\section{The quasi-stationary distribution}
\label{dist:sec:qsd}

The BDC process is absorbed in the end, either at the rupture state $H$ or by
death at $0$; the theory of what a process looks like before that happens is
quasi-stationarity \cite{van2011quasi}. A realisation which has, by chance,
not yet reached fixation has had to hold a middle course. Fly too low and
depletion to zero becomes more and more probable; fly high, a population
booming in number, and the catastrophe clock has all the more to ring
against. What survives hovers, on Daedalus's instruction exactly: neither so
low that the spray soaks the wings, nor so high that the sun reaches them.

A rogue process still going in defiance of the odds is described by the
quasi-stationary distribution: a set of limiting conditional probabilities,
each qualified by non-fixation. The general definition, the Yaglom limit and
the classical existence theory are in
\mextref{m:def:qsd}{the quasi-stationarity definition in \ChM} and
\mextref{m:sec:condmeans}{its treatment of conditional means}, resting on the
classification of birth--death
processes \cite{karlin1957classification}. Here the limit is quoted from
nowhere: it is read off \cref{dist:thm:geometric} by letting $P$ reach its
endpoint.

\subsection{Conditioning on survival}
\label{dist:sec:conditioning}

Let $F(t)$ be the probability of fixation by time $t$, whether by death or by
collapse. The decomposition of \cref{dist:sec:moments-recap} settles it
immediately: a realisation is fixed by $t$ precisely when it is not
productively infected, so
\begin{equation}
F(t)=1-\Ifix(t)=1-I(t)+D(t),
\label{dist:eq:fixprob}
\end{equation}
which reduces to $1-I(t)$ when $\mu=0$. Conditional state probabilities
decompose as
\[
\pr{\Xt=i}=\pr{\Xt=i\mid\text{fixation}}F(t)
+\pr{\Xt=i\mid\text{non-fixation}}\lrb{1-F(t)},
\]
and for $i\notin\{0,H\}$ the process has not been absorbed, so the first term
vanishes and
\begin{equation}
\pr{\Xt=i\mid \text{non-fixation}}
=\frac{p_i(t)}{1-F(t)}=\frac{p_i(t)}{\Ifix(t)}.
\label{dist:eq:condp}
\end{equation}
The quasi-stationary distribution is the limiting form of these:
\begin{equation}
q_i:=\lim_{t\to\infty}\pr{\Xt=i\mid\Xt\notin\{0,H\}}
=\lim_{t\to\infty}\frac{p_i(t)}{\Ifix(t)},\qquad i\ge1.
\label{dist:eq:qsddef}
\end{equation}

\subsection{The limit is geometric}
\label{dist:sec:qsd-geometric}

\begin{theorem}[Quasi-stationary distribution of the BDC process]
\label{dist:thm:qsd}
For all $\mu\ge0$ the quasi-stationary distribution exists and is geometric
on $\{1,2,\ldots\}$ with ratio $1/a$:
\begin{equation}
q_n=\Bigl(1-\frac{1}{a}\Bigr)\Bigl(\frac{1}{a}\Bigr)^{n-1}
=(a-1)\,a^{-n},\qquad n\ge1.
\label{dist:eq:qsd}
\end{equation}
\end{theorem}

\begin{proof}
By \cref{dist:eq:condp} and \cref{dist:cor:sum}, the conditional law at time
$t$ is
\begin{equation}
\pr{\Xt=n\mid\Xt\notin\{0,H\}}
=\frac{p_n(t)}{\sum_{m\ge1}p_m(t)}
=\lrb{1-P(t)}P(t)^{n-1},
\label{dist:eq:condgeom}
\end{equation}
for every $t$: no limit has been taken, and the conditional law is already
geometric. By \cref{dist:thm:geometric}, $P(t)\uparrow1/a$, so each
coordinate converges and the limit is \cref{dist:eq:qsd}.
\end{proof}

The theorem is a corollary of \cref{dist:thm:geometric} in substance as well
as in proof: the quasi-stationary law is not a new object but the endpoint of
a family that was geometric all along. Its moments follow immediately:
\begin{equation}
\langle X\rangle_{\rm QS}=\frac{a}{a-1}=\frac{a}{A},
\qquad
\langle X^2\rangle_{\rm QS}=\frac{a(a+1)}{(a-1)^2},
\qquad
\var{X}_{\rm QS}=\frac{a}{(a-1)^2}=\frac{a}{A^2}.
\label{dist:eq:qsmoments}
\end{equation}
The quasi-stationary mean is exactly the conditional mean burst size
$\E{\Kb\mid\text{burst}}$ of \cref{dist:eq:vinf}, and at $\mu=0$ both also
equal $V_\infty$. Once $\mu>0$ it is $V_\infty$ that falls away from the other
two: at $(1,0.2,0.05)$ it is $13.99$ against
$\langle X\rangle_{\rm QS}=17.23$, because $V_\infty$ averages in the cells
that cleared internally and released nothing. \Cref{dist:fig:qsmeans} traces
all three across $\mu$. That the first two agree is not a coincidence of first
moments: \cref{dist:sec:burst-size} shows the entire burst-size distribution
\emph{is} the quasi-stationary distribution.

\begin{figure}[tbp]
\centering
\includegraphics[width=0.68\textwidth]{fig04_qs_vs_release_means/fig04}
\caption{Quasi-stationary and release means across internal death rates at
$\lambda=1$ and $\delta=0.05$. The quasi-stationary mean
$\langle X\rangle_{\rm QS}=a/(a-1)$ is identically the conditional burst mean
$V_\infty/(1-b)$. At $\mu=0$, $b=0$ and both also equal $V_\infty$; once
$\mu>0$, internally cleared cells contribute zero release and place
$V_\infty$ strictly below the other two; the arrow marks that gap,
$b\langle X\rangle_{\rm QS}=3.25$, at the working point. In the $\mu=0$ case the conditional
mean and the cumulative release share a limit,
$\lim_t J(t)/I(t)=\lim_t V(t)=1+\lambda/\delta$, which diverges as $\delta$
shrinks.}
\label{dist:fig:qsmeans}
\end{figure}

\subsection{The decay rate, and what it validates}
\label{dist:sec:decay-rate}

Quasi-stationarity has a spectral reading, and carrying it out closes a
question left open in \ChM. Restricting the generator to the positive states
gives the killed subgenerator $\mathcal Q$, with
$\mathcal Q_{i,i+1}=\lambda i$, $\mathcal Q_{i,i-1}=\mu i$ for $i\ge2$, and
$\mathcal Q_{i,i}=-(\lambda+\mu+\delta)i$; its missing row mass is the
instantaneous absorption rate, $\delta i$ for $i\ge2$ and $\mu+\delta$ for
$i=1$. Writing $T_{\rm fix}$ for the first of internal extinction and rupture
and summing the positive-state forward equations,
\begin{equation}
\ddt{}\pr{T_{\rm fix}>t}
= -\mu\pr{\Xt=1} - \sum_{i\ge1}\delta i\,\pr{\Xt=i},
\label{dist:eq:survivalloss}
\end{equation}
all probabilities on the right understood before absorption. Ordinary death
drains survival mass from state $1$ only; catastrophe drains it from every
positive state.

\begin{proposition}[Quasi-stationary decay rate]
\label{dist:prop:decay}
If $\nu$ is a quasi-stationary distribution on the positive states and the
sums below are finite, then for some $\vartheta>0$
\begin{equation}
\nu\mathcal Q=-\vartheta\nu,
\qquad
\mathbb P_\nu\{T_{\rm fix}>t\}=\ee^{-\vartheta t},
\qquad
\vartheta=\mu\nu_1+\sum_{i\ge1}\delta i\,\nu_i .
\label{dist:eq:decayrate}
\end{equation}
\end{proposition}

\begin{proof}
Quasi-stationarity writes the unnormalised positive-state law at time $t$ as
$\alpha(t)\nu$; the semigroup property forces
$\alpha(s+t)=\alpha(s)\alpha(t)$, and right-continuity gives
$\alpha(t)=\ee^{-\vartheta t}$. Differentiating at zero yields the
left-eigenvector relation, and summing its coordinates against
\cref{dist:eq:survivalloss} identifies the decay rate.
\end{proof}

The state-dependence matters. Were the catastrophe rate a constant,
independent of the load, its clock would run independently of the
birth--death process, the killed semigroup would factorise into that clock's
exponential times its birth--death counterpart, and the extra clock would
shift the decay rate by that constant while leaving the conditional law
untouched. A rate proportional to the load weights paths by their accumulated
exposure and generally changes the quasi-stationary law itself. It is
therefore no formality that a quasi-stationary distribution survives
conditioning on \emph{neither} extinction nor rupture; \ChM{} leaves that
open, and \cref{dist:thm:qsd} settles it in closed form.

Substituting that law, $\nu_n=(a-1)a^{-n}$, into \cref{dist:eq:decayrate}
gives a check on everything together. Using $\nu_1=(a-1)/a$,
$\sum_i i\nu_i=a/(a-1)$, and then $\mu=\lambda ab$ and
$\delta=\lambda(a-1)(1-b)$ from \cref{dist:eq:ab},
\begin{equation}
\vartheta=\mu\frac{a-1}{a}+\delta\frac{a}{a-1}
=\lambda ab\,\frac{a-1}{a}+\lambda(a-1)(1-b)\,\frac{a}{a-1}
=\lambda b(a-1)+\lambda a(1-b)=\lambda(a-b)=\theta .
\label{dist:eq:decaycheck}
\end{equation}
The eigenvalue of the killed subgenerator is the time constant that has
appeared in every formula since \cref{dist:eq:abtheta}. One line validates the
quasi-stationary law, the decay-rate formula and $AB=\delta/\lambda$ together.

\subsection{Mean productive lifetime}
\label{dist:sec:lifetime}

One integral of $\Ifix$ is a result in its own right, and is also the first
item of the consignment named in \cref{dist:sec:exports}. The expected
duration of productive infection is
\begin{equation}
\E{T_{\rm prod}}=\int_0^\infty\Ifix(t)\,\dd t,
\label{dist:eq:tprodint}
\end{equation}
since $\Ifix(t)=\pr{T_{\rm prod}>t}$.

\begin{proposition}[Mean productive lifetime]
\label{dist:prop:lifetime}
\begin{equation}
\E{T_{\rm prod}}
=\frac{1}{\lambda}\log\!\Bigl(\frac{a}{a-1}\Bigr)
=\frac{1}{\lambda}\log\langle X\rangle_{\rm QS}.
\label{dist:eq:tprod}
\end{equation}
\end{proposition}

\begin{proof}
Write $\Ifix=(I-b)-(D-b)$. From \cref{dist:eq:idihat},
\[
I(t)-b=\frac{B(a-b)}{B+Aw},
\qquad
D(t)-b=-\frac{b(a-b)}{aw-b},
\]
so that, with $w=\ee^{\theta t}$ and $\dd t=\dd w/(\theta w)$,
\[
\int_0^\infty\Ifix\,\dd t
=\frac{a-b}{\theta}\int_1^\infty
\lrb{\frac{B}{w(B+Aw)}+\frac{b}{w(aw-b)}}\,\dd w.
\]
By partial fractions $\frac{B}{w(B+Aw)}=\frac{1}{w}-\frac{A}{B+Aw}$ and
$\frac{b}{w(aw-b)}=-\frac{1}{w}+\frac{a}{aw-b}$, hence
\begin{align*}
\int_1^\infty\frac{B\,\dd w}{w(B+Aw)}
&=\Bigl[\log\frac{w}{B+Aw}\Bigr]_1^\infty
=\log\frac{a-b}{a-1},\\
\int_1^\infty\frac{b\,\dd w}{w(aw-b)}
&=\Bigl[\log\frac{aw-b}{w}\Bigr]_1^\infty
=\log\frac{a}{a-b}.
\end{align*}
Summing, and using $\theta=\lambda(a-b)$,
\[
\int_0^\infty\Ifix\,\dd t
=\frac{a-b}{\lambda(a-b)}
\lrb{\log\frac{a-b}{a-1}+\log\frac{a}{a-b}}
=\frac{1}{\lambda}\log\frac{a}{a-1}.\qedhere
\]
\end{proof}

When $\mu=0$ this is $\lambda^{-1}\log(1+\lambda/\delta)$. The form of the
answer interprets itself. Here $\log\langle X\rangle_{\rm QS}/\lambda$ is the
time that exponential growth at rate $\lambda$ would take to reach the
quasi-stationary level --- the productive lifetime is precisely the time
needed to climb to the height at which the process can hover.
\Cref{dist:fig:lifetime} maps it across the rate plane; at the working
parameters $\E{T_{\rm prod}}=2.847$, implying a matched
constant-rate death rate $d_{\Icell}=0.351$.

That is the cell while it lasts. What it releases when it stops is a separate
question, and the law answering it has already appeared in this section.

\begin{figure}[tbp]
\centering
\includegraphics[width=0.92\textwidth]{fig05_productive_lifetime/fig05}
\caption{Mean productive lifetime
$\E{T_{\rm prod}}=\lambda^{-1}\log\bigl(a/(a-1)\bigr)$. (a) Its dependence on
$\mu\in[0,1]$ and $\delta\in[0.005,1]$ at $\lambda=1$; contours give lifetime
values and the star marks the working point. (b) The productive-survival
curve $\Ifix(t)$ at $(\lambda,\mu,\delta)=(1,0.2,0.05)$, with its integral
$\E{T_{\rm prod}}=2.847$ shaded. The matched classical death rate for \ChPop{}
is $d_{\Icell}=1/\E{T_{\rm prod}}=0.351$.}
\label{dist:fig:lifetime}
\end{figure}
